\manualchapter{Control reference}{sec:controls}
Panel map: \textbf{1} five waveform editors;
\textbf{2} pitch, period and pitch/FM inputs;
\textbf{3} FM, noise mode and level;
\textbf{4} pulse width and Wave Morph;
\textbf{5} audio outputs.

Five editors run down the left edge. To their right, the voice columns run
left to right: Pulse 1, Pulse 2, Wave and Noise. Tuning and modulation are
above level; pulse duty and Wave Morph are below level, above the outputs.

\begin{tabularx}{\textwidth}{@{}l l X@{}}
\toprule Control & Range; default & CV behavior \\\midrule
Tone frequency & $\pm2.5$ oct.; 0 & C4 reference; V/OCT adds octaves. \\
Tone FM depth & $-1$--$+1$; 0 & Adds $aV/5$ octaves; first FM input normals to 5\,V. \\
Pulse duty & 0--3; 2 & Adds $4V/7$, then truncates/clamps. Widths: 12.5, 25, 50, 75\%. \\
Wave morph & 0--5; 0 & Adds $5V/7$, then clamps to 1--5. Settings 0--1 select wave 1. \\
Noise period & 0--7; 0 & Adds $V/2$, floors/clamps, then reverses the chip register. \\
Pulse/noise level & 0--15; 10 & Multiplies by $V/10$; integer steps. \\
Wave level & 0--3; 3 & Off, 25\%, 50\%, 100\% of the wave channel. \\
\bottomrule
\end{tabularx}

Level CV normals to 10\,V at the first column and forwards through the other
columns. Pitch and FM normal through the three pitched voices; pulse-duty
CV normals from Pulse 1 to Pulse 2, then into the wave-morph CV. Insert a
cable to override a normalled source. Wave Morph has no separate depth knob.

The LFSR switch changes the noise sequence. Its CV switches at 2\,V high and
0.01\,V low relative to the panel state. The legacy drawing shows an unused noise SYNC socket; no such input is
provided by the current widget.
changing the LFSR sequence is not a sync trigger. Pitch/FM run each sample;
the other controls update every 16 host samples.
